\phant{KẾT LUẬN VÀ HƯỚNG PHÁT TRIỂN}\label{sec:ketluan}

\subsection{Kết luận}

Nghiên cứu của Diehl, Nader và Gupta đánh giá bốn LLM hàng đầu trên năm biến thể
lập trình song song cho bài toán Mandelbrot. Kết quả cho thấy một bức tranh vừa
lạc quan vừa thận trọng.

Ở mặt tích cực, LLM đã đủ khả năng tạo ra bản nháp mã song song C++ có cấu trúc
đúng. Ba trong bốn mô hình biên dịch 100\% chương trình; sau khi sửa, mọi mã đều
chạy và mở rộng được. ChatGPT-5 đạt 100\% ở cả ba tiêu chí biên dịch, thực thi và
tính đúng — một kết quả đáng chú ý với một bài toán đa mô hình lập trình. Trong
một số trường hợp, mã LLM còn vượt mã con người: bất đồng bộ của ChatGPT-4 và
OpenMP của ChatGPT-5 (theo phần văn bản bài báo gốc; xem lưu ý đối chiếu đồ thị
tại Phần~\ref{sec:ketqua}).

Ở mặt hạn chế, khoảng cách giữa “biên dịch được” và “chạy đúng” rất lớn với các
mô hình yếu. ChatGPT-4 biên dịch 100\% nhưng chỉ 20\% (1/5) chương trình chạy
không lỗi; Claude đạt 40\% theo Bảng~III gốc (3/5 = 60\% theo Bảng~I--II); LLaMA
đạt 0\% theo các bảng tổng hợp. Các lỗi tập trung ở ba nhóm: kiểm soát số luồng (nhầm OpenMP với TBB), vòng đời coroutine (\code{final\_suspend},
\code{promise\_type}, giải phóng hai lần) và định dạng đầu ra (P4/P5). Đây đều là
lỗi ngữ nghĩa hoặc lỗi hệ sinh thái, không phải lỗi cú pháp.

Kết luận trung tâm là: năng lực của LLM trong HPC không nằm ở việc viết đúng cú
pháp, mà ở việc hiểu đúng \emph{ngữ cảnh thực thi} — backend nào, phiên bản nào,
tài nguyên nào, chi phí nào. Khoảng cách lớn nhất so với chuyên gia không phải ở
số dòng mã hay độ phức tạp thuật toán, mà ở hiểu biết về vòng đời tài nguyên và
mô hình truyền thông. Đây là lý do LLM có thể rất hữu ích như trợ lý, nhưng chưa
thể thay thế kỹ sư HPC.

\subsection{Khuyến nghị}

Từ kết quả thực nghiệm, có thể rút ra năm khuyến nghị thực dụng:

\begin{enumerate}
  \item \textbf{Coi mã LLM là bản nháp, không phải sản phẩm.} Luôn biên dịch với
        cảnh báo bật, chạy trên dữ liệu nhỏ và kiểm tra đầu ra trước khi tin
        tưởng.
  \item \textbf{Đo mở rộng thay vì tin vào cờ dòng lệnh.} Kiểm tra backend thực
        tế và xác nhận số luồng có hiệu lực bằng phép đo, không chỉ bằng việc đọc
        mã.
  \item \textbf{Soát kỹ coroutine và MPI.} Với coroutine, kiểm tra
        \code{final\_suspend}, \code{promise\_type} và quyền sở hữu frame. Với
        MPI, kiểm tra kích thước buffer, kiểu dữ liệu, trường hợp một tiến trình
        và mốc đo thời gian.
  \item \textbf{Kiểm tra định dạng đầu ra.} Lỗi P4/P5 có thể tạo ảnh đen mà không
        báo lỗi; cần kiểm tra header, độ sâu bit và kích thước tệp.
  \item \textbf{Viết prompt nêu rõ ràng buộc phi chức năng.} Số luồng/phân vùng
        cấu hình được, đo thời gian loại trừ I/O, định dạng đầu ra, phiên bản
        chuẩn C++, và — với MPI — yêu cầu chia việc thực sự thay vì tính lặp.
\end{enumerate}

Một khuyến nghị ở mức tổ chức: nên xem LLM như một thành viên nhóm cần được
review, không phải một công cụ tự động. Quy trình review cần có danh mục kiểm tra
riêng cho mã song song, vì các tiêu chí thông thường (biên dịch được, qua unit
test) không đủ phát hiện lỗi tôn trọng tài nguyên và lỗi đo đạc.

\subsection{Hướng phát triển}

Bài báo đề xuất nhiều hướng mở rộng, có thể nhóm thành bốn trục.

\subsubsection{Mở rộng bài toán}

Cần đánh giá trên các benchmark truyền thông dày: stencil có trao đổi halo, FFT
song song, bộ giải lặp dùng collective MPI. Đây là nơi lỗi đồng bộ và chi phí
truyền thông thực sự lộ diện. Bên cạnh đó, các ứng dụng khoa học phức tạp như
phần tử hữu hạn hoặc sai phân hữu hạn sẽ kiểm tra khả năng sinh mã cho bài toán
có cấu trúc dữ liệu và phụ thuộc phức tạp.

\subsubsection{Mở rộng mốc so sánh}

Cần so sánh với các phiên bản mã con người được tối ưu cao hơn, không chỉ mã sách
giáo khoa. Đồng thời, nên dùng công cụ profiling và gỡ lỗi để giải thích vì sao
một số mã LLM nhanh hơn hẳn các mã khác — điều mà phép đo thông lượng đơn thuần
không trả lời được.

\subsubsection{Mở rộng mô hình}

Cần bổ sung các mô hình mới như Gemini, DeepSeek-Coder, Grok và các mô hình chuyên
mã. Nên đánh giá nhiều lần sinh cho mỗi cấu hình để có phân tích thống kê, và thử
nghiệm với nhiều cách diễn đạt prompt, nhiệt độ khác nhau. Đây là điều kiện để
chuyển từ nghiên cứu tình huống sang đánh giá có thể khái quát hóa.

Cách bổ sung cụ thể (chọn mô hình, giữ nguyên prompt/toolchain, lặp lại phép đo và
đối chiếu với nghiên cứu dùng phương pháp khác) cùng bối cảnh từ ba công trình độc
lập cuối năm 2026 \cite{bantel2026agentic,cui2026hpcbench,nogueira2026openmp} đã
được trình bày tại Mục~\ref{subsec:mo-hinh-hien-nay}. Các kết quả độc lập này bước
đầu cho thấy thế hệ mô hình mới (kể cả trọng số mở) cải thiện rõ tính đúng và tăng
tốc trên OpenMP/tasking, nhưng vẫn gặp giới hạn khi bị tuần tự hóa ngầm hoặc siết
chặt ngưỡng speedup; để kiểm chứng trên cùng thước đo với coroutine C++20 và
MPI+OpenMP, vẫn cần chạy thực nghiệm trực tiếp với bộ prompt P1/P2/P3 của bài toán
Mandelbrot.

\subsubsection{Mở rộng tiêu chí}

Nên so sánh LLM với các công cụ truyền thống như Intel Parallel Advisor, công cụ
phân tích tĩnh và công cụ profiling. Điều này cho thấy LLM bổ sung hay thay thế
những công cụ nào trong quy trình phát triển. Cuối cùng, cần xây dựng bộ tiêu chí
đánh giá mã song song do AI sinh, bao gồm tính đúng, khả năng mở rộng, tính tái
lập và chi phí bảo trì — chứ không chỉ dừng ở tỉ lệ biên dịch.

\subsection{Kết luận cuối}

Bài báo đặt nền móng cho một hướng nghiên cứu còn non trẻ: đánh giá LLM trong bối
cảnh HPC thực thụ. Giá trị lớn nhất của nó không nằm ở việc tuyên bố mô hình nào
“tốt nhất”, mà ở việc chỉ ra \emph{loại lỗi nào} xuất hiện, \emph{vì sao} chúng
xuất hiện, và \emph{cần đo gì} để phát hiện. Đó là nền tảng để xây dựng các quy
trình đánh giá nghiêm túc hơn, các mô hình chuyên biệt hơn, và các công cụ hỗ trợ
lập trình HPC đáng tin cậy hơn trong tương lai.
